\section{Repeatable observations and causal realization}\label{sec:causal59}
Terminal numerical accuracy and accuracy of a generated observation process
are different requirements. Repeating a nondisturbing observation can
separate arbitrarily close, but distinct, numerical responses. We now
formulate this distinction within the same finite-label resource model.

Let $H$ be a compact Hausdorff group, let $\mathcal A$ be a finite nonempty
command alphabet generating a dense subgroup, and let $\mathcal S$ and
$\mathcal J$ be finite nonempty sets. For each $s,j$ fix a continuous law
\[
 p_{s,j}:H\longrightarrow\Delta_{M_j}.
\]
The physical state starts at $g=e$. A command $a$ replaces $g$ by $u_ag$
and emits a fixed null symbol. A probe $j$ leaves $g$ unchanged and emits
a fresh draw from $p_{s,j}(g)$, conditionally on the public past
and current physical state. At each step the next command or probe may be chosen from the entire
observed transcript by a randomized policy. The seed is specified at
initialization. Probes are repeatable statistical observations; when
$p_{s,j}$ is obtained from a density matrix, this is an oracle model, not
repeated nondisturbing measurement of one unknown quantum system.

\begin{definition}[Causal finite-label machine]\label{def:causal59}
At horizon $N$ a machine has registers $S_0,\ldots,S_N$, initialization
rows $E_s$, and, for each control $c$ and possible output $y$, a
nonnegative matrix $K_{t,c,y}:S_{t-1}\to S_t$ such that
$\sum_yK_{t,c,y}$ is row-stochastic. Its entries jointly specify the
emission and next label. The matrices may depend on $t,N$ but on the past
only through the retained label and current control. Its error is
\[
 \sup_{s,\pi}\TV\bigl(\mathbb P^{\pi}_{s,N},
                            \mathbb Q^{\pi}_{s,N}\bigr),
\]
where the laws include the complete control--output transcript under the
same policy $\pi$. Write $\mathsf{PW}_{N,\varepsilon}$ for the minimum
peak register size at error at most $\varepsilon$.
\end{definition}
Every advertised label and all retained private randomness are counted.
The horizon, external clock and arbitrary real tables remain free. The
policy is an external experimenter, not an uncharged private memory for
the machine. Outputs may influence future controls, but are not supplied
again as a free history input. A sampled output register, when charged,
adds the separate fixed output-alphabet cost. The error above is one
budget for the entire transcript, not a per-step error budget.

For $(s,g)$ define its future response function
\begin{equation}\label{eq:futuretype59}
 \Theta_{s,g}(x,j)=p_{s,j}(xg),\qquad x\in H,
 \qquad \mathcal X=\{\Theta_{s,g}:s\in\mathcal S,\ g\in H\}.
\end{equation}
Equip the finite product of continuous-function spaces with the uniform
norm. Put $n=|\mathcal X|$ when this set is finite, and $n=\infty$
otherwise. This is a behavioral quotient, not the dimension of a linear
predictive representation. In particular, infinitely many distinct
predictive rows do not by themselves imply infinite nonnegative
realization size for a general hidden Markov process
\cite[Theorem~2 and Section~4]{SJR2004}. Repeatable observations are the
additional input in the theorem below.

\begin{theorem}[Uniform process error and finite physical actions]
\label{thm:causalclassification59}
For every fixed $0\le\varepsilon<1$ the following are equivalent:
\begin{enumerate}
\item $\sup_N\mathsf{PW}_{N,\varepsilon}<\infty$;
\item $n<\infty$;
\item some open normal subgroup $U\lhd H$ makes every $p_{s,j}$ constant
      on each coset of $U$;
\item a stationary exact causal realization exists with permutation
      command updates and state-preserving probe emissions.
\end{enumerate}
If $n<\infty$ and $0\le\varepsilon<1/2$, then
\begin{equation}\label{eq:causalminimum59}
                   \mathsf{PW}_{N,\varepsilon}=n
                   \quad\hbox{for all sufficiently large }N.
\end{equation}
If $n=\infty$, then, for every $\varepsilon<1$ and every $k\ge1$,
there is $L_{k,\varepsilon}<\infty$ such that every admissible realization
satisfies
\begin{equation}\label{eq:causaloccupation59}
 \#\{t\in\{1,\ldots,N\}:|S_t|\le k\}\le L_{k,\varepsilon}.
\end{equation}
The same occupation conclusion holds for finite $n$, $\varepsilon<1/2$
and $k<n$. Intervening registers are unrestricted. No return condition
is imposed on the command alphabet: the probes themselves leave the
physical state unchanged and supply the padding used in the proof.
\end{theorem}
The threshold $1$ in the boundedness criterion is sharp: total variation
never exceeds one, so at tolerance one a single label always suffices.
The elementary one-continuation proof of \eqref{eq:causalminimum59}
uses $\varepsilon<1/2$. Theorem~\ref{thm:causalstrong62} below
proves the full minimum for every $\varepsilon<1$ by a different
tail-event argument; the present proof and its stronger one-cut test
remain useful in their stated range.

\begin{lemma}[A common separating itinerary]\label{lem:itinerary59}
Choose $m\ge2$ different functions in \eqref{eq:futuretype59}, represented
by executable seed--prefix pairs. For every $\delta>0$ there is one
finite deterministic control itinerary with disjoint output events
$A_1,\ldots,A_m$ such that its target law from the $i$th state assigns
probability at least $1-\delta$ to $A_i$.
\end{lemma}
\begin{proof}
Equality of two future response functions is preserved and reflected by
left multiplication of their physical states by any $u\in H$: replace
$x$ in \eqref{eq:futuretype59} by $xu$, a bijection of $H$. Positive
executable products are dense in $H$. Indeed positive powers approximate
the inverse of every generator in a compact group, so the closed positive
semigroup contains the dense generated group.

List the $\binom m2$ pairs. At each stage the current two states for that
pair still have different future response functions. By continuity and
positive-word density, a finite executable word followed by one probe
separates some coordinate of their output laws. Append that word and
record the coordinate, then proceed to the next pair. This constructs
one itinerary, not a family of incompatible pair-dependent suffixes.
No reset or inverse word is required. Let $J=\binom m2$, and let
$\mu_{i,r}$ be the selected coordinate probability at the $r$th probe
location when the initial state is the $i$th representative. There is
$\Delta>0$ such that every pair of rows of this finite matrix differs
by at least $\Delta$ in some coordinate.

At each location repeat the chosen probe $L$ times. These repetitions
leave the physical state fixed. The empirical coordinate frequencies
$\widehat\mu_r$ obey
\[
 \mathbb P_i\{\max_r|\widehat\mu_r-\mu_{i,r}|\ge\Delta/3\}
       \le 2J\exp(-2L\Delta^2/9).
\]
For completeness, the scalar bound follows from
$\E e^{z(B-\E B)}\le e^{z^2/8}$ for a Bernoulli variable, exponential
Markov inequality, and optimization in $z$; a union bound gives the
display. The moment bound follows by differentiating the logarithmic
moment generating function twice: its second derivative is the variance
of a tilted Bernoulli variable and is at most $1/4$.
Choose $L\ge1$ with the displayed upper bound at most $\delta$.
The events
$A_i=\{\max_r|\widehat\mu_r-\mu_{i,r}|<\Delta/3\}$ are disjoint,
because $2\Delta/3<\Delta$, and have the required probabilities.
\end{proof}

\begin{lemma}[Testing across a finite register]\label{lem:testingcut59}
Suppose $m$ equal-length seed--prefix pairs at cut $t$ admit an itinerary
of length $\ell$ as in Lemma~\ref{lem:itinerary59}. For every
$N\ge t+\ell$, an error-$\varepsilon$ causal realization satisfies
\begin{equation}\label{eq:testingcut59}
                  |S_t|\ge m(1-\varepsilon-\delta).
\end{equation}
If $\varepsilon+\delta<1/2$, the stronger bound $|S_t|\ge m$ holds.
\end{lemma}
\begin{proof}
Use the same deterministic suffix for all representatives and pad its
end with probes if necessary. Marginalize all prefix outputs. Write
$\alpha_i(z)$ for the resulting cut-label distribution and $Q_z$ for
the suffix-output law conditional on label $z$. Causality makes $Q_z$
independent of the chosen prefix. Total variation contracts under
marginalization, so correctness gives
\[
 1-\delta-\varepsilon
     \le\sum_z\alpha_i(z)Q_z(A_i).
\]
This uses no conditional correctness assertion after a rare observed
prefix. Summing in $i$, using $\alpha_i(z)\le1$, and then disjointness,
yields
\[
 m(1-\delta-\varepsilon)
 \le\sum_z\sum_iQ_z(A_i)\le |S_t|.
\]
If $1-\delta-\varepsilon>1/2$, each $i$ has a positive-mass label $z_i$
with $Q_{z_i}(A_i)>1/2$. Disjointness forces the labels $z_i$ to be
different. This proves the stronger assertion.
\end{proof}

\begin{proof}[Proof of Theorem~\ref{thm:causalclassification59}]
The map $g\mapsto\Theta_{s,g}$ is continuous in the uniform norm.
This follows from continuity on the compact product $H\times H$ and a
finite-cover argument, uniformly in the first coordinate. Thus each
seed's image is compact. Executable images are dense in it. If
$\mathcal X$ is infinite, arbitrarily many distinct types therefore
have executable representatives; if it is finite, every type has one.

Choose finitely many such representatives and pad their prefixes to a
common length $t_0$ by any fixed probe. Their physical states are
unchanged. Apply Lemma~\ref{lem:itinerary59}. For every
$t_0\le t\le N-\ell$, pad further with the same probe, marginalize its
outputs, and apply Lemma~\ref{lem:testingcut59}. The itinerary and
$\ell$ do not depend on $t,N$ or on the machine.

If $n=\infty$, choose $m$ with $m(1-\varepsilon)>k$ and then choose
$\delta>0$ with $m(1-\varepsilon-\delta)>k$. Every interior cut has
width greater than $k$. At most $t_0+\ell$ positive cuts lie outside
this interval, including when it is empty. This proves
\eqref{eq:causaloccupation59}, divergence, and the implication (i)$\Rightarrow$(ii).
For finite $n$ and $\varepsilon<1/2$, use $m=n$ and
$0<\delta<1/2-\varepsilon$ when $n\ge2$. All sufficiently interior
cuts require $n$ labels. The case $n=1$ follows from nonempty registers.

For finite $\mathcal X$, the maps
\[
 \Theta(x,j)\longmapsto\Theta(xu,j),\qquad u\in H,
\]
define a continuous permutation action of $H$ on $\mathcal X$.
Continuity follows from uniform continuity of each of the finitely many
functions. Its kernel $U$ is open and normal. Evaluating at $x=e$
shows that every $p_{s,j}$ is constant on $U$-cosets, proving
(ii)$\Rightarrow$(iii). Conversely (iii) gives at most
$|\mathcal S|[H:U]$ response types. More economically, use the $n$
distinct types as labels, initialize at $\Theta_{s,e}$, update by the
permutation above, and emit $\Theta(e,j)$ on probe $j$ without changing
the label. This produces the exact conditional response at each step.
Induction on transcript length proves equality for every adaptive
policy, hence (iv) and (i). It also supplies the upper bound $n$ in
\eqref{eq:causalminimum59}. The occupation assertion for $k<n$ follows
from the stronger testing bound just established.
\end{proof}

\begin{corollary}[Connected groups]\label{cor:connectedprocess59}
If $H$ is connected, uniformly bounded causal width at any fixed total
transcript error below one is possible exactly when every $p_{s,j}$ is
constant on $H$. If at least one is nonconstant, every fixed width has
the occupation bound \eqref{eq:causaloccupation59}.
\end{corollary}
\begin{proof}
A continuous action of a connected group on a finite discrete set is
trivial. The assertion follows from the theorem.
\end{proof}

\begin{remark}[Why terminal purification is insufficient]\label{rem:correlation59}
A two-label variable chosen once with equal masses and then repeatedly
reported has the correct Bernoulli$(1/2)$ marginal at every time, but its
length-$L$ law is supported on the two constant strings. Its total
variation distance from $L$ independent fair bits is $1-2^{1-L}$.
Thus even exact numerical marginals do not supply a path-law realization.
The causal theorem concerns an actual nonnegative emission instrument,
not a list of separately correct decoders. It neither classifies
arbitrary irreversible hidden-state systems nor asserts a measurement
procedure for noncommuting quantum observables.
\end{remark}
